\section{B. 代数与几何}

\subsection{B1. 从矩阵到群论}

\subsubsection{群的基本概念}

\textbf{群 $(G, \cdot)$：} 集合 + 运算，满足封闭性、结合律、单位元、逆元。

\textbf{关键群：}
\begin{center}
\begin{tabular}{llll}
\toprule
\textbf{群} & \textbf{定义} & \textbf{维度} & \textbf{应用} \\
\midrule
$SO(2)$ & 2D 旋转 & 1 & 2D 旋转 \\
$SO(3)$ & 3D 旋转 & 3 & 刚体旋转 \\
$SE(3)$ & 3D 刚体运动 & 6 & 机器人学、动画 \\
$U(1)$ & 复数单位圆 & 1 & 电磁规范对称 \\
$SU(2)$ & $2\times2$ 酉矩阵，$\det=1$ & 3 & 自旋、四元数 \\
$SU(3)$ & $3\times3$ 酉矩阵，$\det=1$ & 8 & 量子色动力学 \\
$O(3)$ & 正交矩阵 & 3 & 包含反射 \\
$GL(n)$ & 可逆矩阵 & $n^2$ & 一般线性变换 \\
$SL(n)$ & $\det=1$ 的矩阵 & $n^2-1$ & 保体积变换 \\
\bottomrule
\end{tabular}
\end{center}

\subsubsection{李群与李代数}
\label{subsec:lie-intro}

\textbf{李群：} 既是群又是光滑流形（群运算光滑）。

\textbf{李代数 $\mathfrak{g}$：} 李群在单位元处的切空间，配备李括号 $[\cdot, \cdot]$。

\textbf{指数映射：} $\exp: \mathfrak{g} \to G$

$\mathfrak{so}(3)$ 的具体形式：
\[
\mathfrak{so}(3) = \{A \in \mathbb{R}^{3\times 3} : A^T = -A\}
\]

反对称矩阵参数化：
\[
[\omega]_\times = \begin{pmatrix} 0 & -\omega_3 & \omega_2 \\ \omega_3 & 0 & -\omega_1 \\ -\omega_2 & \omega_1 & 0 \end{pmatrix}
\]

Rodrigues 公式：
\[
R = \exp([\omega]_\times) = I + \frac{\sin\theta}{\theta}[\omega]_\times + \frac{1-\cos\theta}{\theta^2}[\omega]_\times^2
\]
其中 $\theta = \|\omega\|$。

\textbf{$SU(2)$ 与四元数：} 四元数 $q = a + bi + cj + dk$，单位四元数构成 $SU(2)$（同构）。$SU(2) \to SO(3)$ 是 2:1 的覆盖映射（旋转 $2\pi$ 后四元数变号，需要 $4\pi$ 才回到自身——自旋 $1/2$ 粒子的数学根源）。

\textbf{经典李群的完整分类见第~\ref{sec:lie-group}~节；李代数结构常数、Killing 形式、Dynkin 图与表示论见第~\ref{sec:lie-algebra}~节及其后。}

\subsubsection{表示论基础}

\textbf{表示：} 群 $G$ 到 $GL(V)$ 的同态 $\rho: G \to GL(V)$。

\textbf{不可约表示：} 没有非平凡不变子空间的表示。

\textbf{在物理中：}
\begin{itemize}
  \item 角动量的量子化来自 $SO(3)$ 的不可约表示（维数 $2l+1$）
  \item 粒子分类来自庞加莱群的不可约表示
  \item 球谐函数 $Y_l^m$ 就是 $SO(3)$ 的第 $l$ 个不可约表示的基函数
\end{itemize}

\textbf{在 CG 中：}
\begin{itemize}
  \item 球谐光照（Spherical Harmonics Lighting）
  \item 等变神经网络（Equivariant Neural Networks）
\end{itemize}

%--------------------------------------------------------------
\subsection{B2. 微分几何}

\subsubsection{曲线与曲面}

\textbf{曲线} $\gamma: [a,b] \to \mathbb{R}^3$：
\begin{itemize}
  \item 弧长参数化：$s = \int_0^t \|\gamma'(\tau)\|\,d\tau$
  \item 曲率：$\kappa = \|\mathbf{T}'(s)\|$
  \item 挠率：$\tau$
  \item Frenet 标架：$(\mathbf{T}, \mathbf{N}, \mathbf{B})$
\end{itemize}

\textbf{曲面} $\mathbf{r}(u,v)$：
\begin{itemize}
  \item 第一基本形式：$I = E\,du^2 + 2F\,du\,dv + G\,dv^2$（度量）
  \item 第二基本形式：$II = L\,du^2 + 2M\,du\,dv + N\,dv^2$（弯曲）
  \item 高斯曲率：$K = \dfrac{LN - M^2}{EG - F^2}$（\textbf{内蕴量}，不依赖嵌入）
  \item 平均曲率：$H = \dfrac{EN - 2FM + GL}{2(EG - F^2)}$
\end{itemize}

\textbf{高斯绝妙定理（Theorema Egregium）：} 高斯曲率是内蕴的，只依赖第一基本形式。

在 CG 中：网格曲率估计、平均曲率流（网格平滑）、高斯曲率（展开/参数化）。

\subsubsection{流形与切空间}

\begin{itemize}
  \item \textbf{流形 $M$：} 局部像 $\mathbb{R}^n$ 的拓扑空间
  \item \textbf{切空间 $T_pM$：} 点 $p$ 处所有“方向导数”构成的向量空间
  \item \textbf{切丛 $TM$：} 所有切空间的并集
  \item \textbf{余切空间 $T_p^*M$：} 切空间的对偶（1-形式生活的地方）
\end{itemize}

\subsubsection{黎曼几何}

\textbf{黎曼度量 $g$：} 每个切空间上的内积，光滑地依赖于点。
\[
g = g_{ij}\,dx^i \otimes dx^j
\]

\textbf{Christoffel 符号：}
\[
\Gamma^k_{ij} = \frac{1}{2}g^{kl}\left(\frac{\partial g_{il}}{\partial x^j} + \frac{\partial g_{jl}}{\partial x^i} - \frac{\partial g_{ij}}{\partial x^l}\right)
\]

\textbf{黎曼曲率张量：}
\[
R^l_{ijk} = \partial_j \Gamma^l_{ik} - \partial_k \Gamma^l_{ij} + \Gamma^m_{ik}\Gamma^l_{mj} - \Gamma^m_{ij}\Gamma^l_{mk}
\]

\textbf{测地线方程：}
\[
\frac{d^2 x^k}{dt^2} + \Gamma^k_{ij}\frac{dx^i}{dt}\frac{dx^j}{dt} = 0
\]

在 CG 中：网格上的测地线计算、测地线距离（形状匹配）、黎曼优化。

\subsubsection{纤维丛（Fiber Bundle）}

\textbf{直觉：} 一个空间，每一点上面“粘着”一个纤维（向量空间或其他结构）。

\textbf{例子：}
\begin{itemize}
  \item 切丛 $TM$：每一点的纤维是切空间
  \item 余切丛 $T^*M$：每一点的纤维是余切空间
  \item 主丛：纤维是群 $G$（规范场论的数学语言）
\end{itemize}

\textbf{在物理中：} 规范场（电磁场、Yang--Mills 场）是主丛上的联络。